\section{Inside the group: size and leader}\label{sec:inside}

Does the effect grow with the number of pilots, and does the pilot ahead fly differently from the one
behind?

\subsection{Size of the group}\label{sec:size}

The group time is cut where the group changes, so every stretch has a fixed number of pilots
$n$ (2, 3, 4, 5--6, 7+), and a flight can contribute to several sizes. Two contrasts of
Eq.~\eqref{eq:contrast} follow the size: $\G_n-\A$ within the flight, and $\G_n-\G_2$ at the same
site-day and hour, which asks whether a group does more than a pair.

\begin{figure}
\centering
\includegraphics[width=\textwidth]{report_size.pdf}
\caption{By the number of pilots $n$, with 5--95\% bands. Green: $\G_n-\A$ within the flight. Grey, open:
$\G_n-\G_2$ at the same site-day and hour. Values in Table~\ref{tab:size-contrasts}.}
\label{fig:size}
\end{figure}

Pooled, $\Hloc(1000~\mathrm{s})$ rises with $n$ while the same flights alone hardly change
(Appendix~\ref{app:curves}; values in Table~\ref{tab:size}). Within the flight the effect grows from
\DHlocThousandFlightTwo{} with 2 pilots to \DHlocThousandFlightSeven{} with 7+ (Figure~\ref{fig:size},
green); there the day and the site are fixed, so the growth comes from the hour or from the flights
that join large groups. At the same site-day and hour groups of 3 to 6 progress like pairs; only 7+
stands apart, 3\% of the group time, perhaps competition gaggles flying a set task, which the archive
does not flag. At the same place and hour larger groups spend more time climbing and glide flatter,
but a little slower, than pairs, with most of the change from 2 to 4 pilots. On six single site-days
the differences between days exceed those between sizes (Appendix~\ref{app:sitedays}).

\subsection{Leader and follower}\label{sec:mechanisms}

In groups of two pilots, the signed median delay $\ell_{\rm med}$ of the pair interval
(Appendix~B of the note) says who passes first: the leader, when the other passes at least 15~s later;
otherwise the two fly abreast (\Abreast\% of the pairs). Leader and follower are compared within the
same group, the group being the stratum of Eq.~\eqref{eq:contrast}
(Appendix~\ref{app:leaders}). The pilot ahead climbs more, climbs a little faster, glides faster and
glides as flat. The share of searches that end in a climb does not differ (\LeadFound{}
points): on this measure the follower shows no sign of finding thermals through the leader. How
groups form and break and whether the same pilots lead are in the note on the structure of groups.

\paragraph{What we keep.} Most of the effect is there with two pilots. Beyond a pair, at the same
site-day and hour, size adds climbing and flatter but slightly slower glides, and no straighter
progress, except in the rare groups of 7 or more. The leader is a little faster, and ends searches in a climb as often
as the follower.
